% problem-000386 4 硅化学与缩聚
\section*{4. 硅化学与缩聚}
\textit{下列为分问。}

\subsection*{4-1}
硅酸根基本结构单元是硅氧四⾯体 $\mathrm { ( S i O _ { 4 } ) }$ ，可以⽤三⻆形表⽰硅氧四⾯体的平⾯投影图形，三⻆形的顶点代表⼀个O原⼦，中⼼是Si和⼀个顶⻆O的重叠。若⼲硅氧四⾯体可以连接成层型结构，如图4-1(a)所⽰，两个硅氧四⾯体共⽤⼀个氧原⼦，⼜可以将其简单地表⽰为图4-1(b)。

每⼀种层型结构可以⽤交于⼀点的�元环及环的数⽬来描述，例如图4-1(b)可以描述为 $6 ^ { 3 }$ 。

图4-1(c)⾄(h)给出在沸⽯分⼦筛中经常⻅到的⼆维三连接层结构，图4-1(c)和图4-1(d)可分别描述为4.82 和(4.6.8)(6.82)。

试将对图4-1(e)⾄(h)的描述分别写答卷上。

（图：）
(a)

（图：）

（图：）
(c)

(b)
（图：）
(d)

（图：）
(e)

（图：）
(f )

（图：）
(g)

（图：）
(h)
图4-1 ⼏种层状硅酸根结构的图⽰⽅法

\subsection*{4-2}
分⼦筛⻣架中存在⼀些特征的笼型结构，⻅图4-2。笼型结构是根据构成它们的多⾯体的各种�元环进⾏描述的。例如图4-2(a)所⽰的笼可以描述为45.52；⽽图4-2(b)的笼可以描述为 $4 ^ { 6 } . 6 ^ { 2 } . 6 _ { \mathrm { ~ o ~ } } ^ { 3 }$

试将对图4-2(c)⾄图4-2(h)的描述写在答卷上。

（图：）
(a)

（图：）
(b)

（图：）
(c)

（图：）
(d)

（图：）
(e)

（图：）
(f )

（图：）
(g)

（图：）
(h)

\subsection*{4-3}
分⼦筛⻣架中的笼型结构可以采⽤下述⽅法画出：将最上⼀层放在平⾯图的最内圈，下⼀层放在次内圈，依此类推，直到最底层放在平⾯图的最外圈。图中各原⼦之间的连接关系、各种环及其位置当然要与笼型结构中保持⼀致。例如图4-3(a)所⽰的笼，可以表⽰为图 $4 \mathrm { - } 3 ( \mathrm { a } ^ { \prime } )$ C

图4-3(b)⾄图4-3(d)结构不完整地图⽰在图4-3(b′)⾄图4-3(d′)，即只给出了结构的最上层和最下层的图${ \overline { { \overline { { \prime } } \mathrm { { I \hat { \ddots } } } } } } ,$ 试在答卷上画出完整的平⾯图。

\subsection*{(5)}
（图：）
(a)

（图：）
(b)

（图：）
(c)

（图：）
(d)

（图：）
(a′)

（图：）
(b′)

（图：）
(c′)

（图：）
(d′)
